%% Supplementary material of CMPB-D-26-02177 (revision R1), assembled 25 Oct 2026.
%% All tables are generated by scripts/make_revision_tables.py (tables/*.tex), all figures by
%% scripts/make_revision_figures.py (figs/figS_*.pdf).  Numbers: tables/numbers.md.
\documentclass[10pt,a4paper]{article}
\usepackage[utf8]{inputenc}
\usepackage[T1]{fontenc}
\usepackage[margin=20mm]{geometry}
\usepackage{amsmath,amssymb,booktabs,graphicx,caption,float,longtable}
\usepackage[numbers]{natbib}
\captionsetup{font=small,labelfont=bf}
\setlength{\emergencystretch}{2em}
\graphicspath{{figs/}}
% tables are input from tables_supp/, copies of tables/ with [t] -> [H] and table* -> table (made by the build step)
\renewcommand{\thetable}{S\arabic{table}}
\renewcommand{\thefigure}{S\arabic{figure}}
% every table carries its number explicitly (S1b, S11a/b, S14b break a running counter)
\newcommand{\stab}[1]{\renewcommand{\thetable}{#1}}
\title{Supplementary material\\[4pt]\large Geodesic Decay Analysis: Quantifying the Spatial Resolution of Intracardiac Pace-Mapping Score Maps on Triangulated Atrial Surfaces}
\author{S.\ Willert, V.\ Maslova, A.\ Zaman, D.\ Frank, E.\ Lian}
\date{Revision R1, October 2026}
\begin{document}
\maketitle
\noindent All tables and figures of this supplement are produced by two scripts of the accompanying repository (\texttt{make\_revision\_tables.py} and \texttt{make\_revision\_figures.py}) from the archived result files of the validation and analysis scripts, which were computed from the meshes and score fields; the per-map values of Table~S12 are also provided as a CSV file.  Unless stated otherwise, ``56 meshes'' are the 56 non-fragmented clinical meshes (40 distinct atrial reconstructions, because maps of different regions of one chamber share the reconstruction) and ``55 maps'' the clinical cohort of the main text.

\section*{S1 Heat Method against analytic geodesics}
Distance errors on spheres and open cylinders with analytic geodesics (Table~S1, Fig.~S1), and the width recovered with each distance definition relative to the analytic distance on the same fields (Table~S1b).
\stab{S1}\input{tables_supp/S1_analytic_surfaces}
\stab{S1b}\input{tables_supp/S1b_width_by_solver}
\begin{figure}[htbp]\centering\includegraphics[width=0.95\linewidth]{figS_analytic_surfaces.pdf}
\caption{Distance error of the Heat Method, the exact solver and Dijkstra's algorithm against analytic geodesics on spheres ($R$ = 25 mm; 642, 2{,}562 and 10{,}242 vertices) and open cylinders ($R$ = 20 mm; regular and jittered triangulations): signed relative error and 95th percentile of the absolute relative error, mean over five sources per surface.}\end{figure}

\section*{S2 Heat Method against exact geodesics on the clinical meshes}
\stab{S2}\input{tables_supp/S2_heat_vs_exact_meshes_long}

\section*{S3 Saturating sensitivity protocol: noise-free peak at the upper score limit}
In this protocol ($A$ = 35, $B$ = 65) the noise-free peak equals the upper score limit, so that noisy fields saturate in a plateau around the source, which the tie rule of the origin exploits.  Tied maxima are common at large widths: in a test on 14 meshes at 30\,mm and 10\,\% noise every field had a tied maximum, and the mean bias with the first tied vertex, the plateau centre and a random tied vertex as origin was $+1.99$, $+0.22$ and $+0.88$\,mm (SD 4.38, 1.19 and 3.23\,mm); the plateau-centre rule was chosen after these simulations.  Noise-free fields and the known-origin mode were evaluated in this protocol.
\stab{S3}\input{tables_supp/S3_validation_grid}
\stab{S3b}\input{tables_supp/S3b_known_origin_wide}

\section*{S4 Saturating sensitivity protocol: widths outside 3--30 mm}
\stab{S4}\input{tables_supp/S4_range}

\section*{S4b Primary protocol: full grid}
The primary protocol of the main text (Table~1, Fig.~4) places the noise-free peak at 90 (amplitude 35 over a baseline of 55, clip at 50 and 100, otherwise as in Table~S3).  The saturating sensitivity protocol places the noise-free peak at 100.  Their comparison captures the combined effects of changing the baseline and of clipping at both score limits (at 10\,\% noise a baseline vertex falls below 50 before clipping with a probability of about 8\,\% instead of $<0.01$\,\%).  Tied maxima occurred in at most 2.1\,\% of the fields up to 30\,mm at 10\,\% noise (4.0 and 9.6\,\% at 40 and 50\,mm).  At 2\,mm and 10\,\% noise one of the 1{,}120 fits ran to the parameter bound (mean bias $+0.22$\,mm, SD 5.9\,mm; without that fit SD 0.49\,mm, central 95\,\% of the estimates 1.5--2.7\,mm).
\stab{S4b}\input{tables_supp/S4b_peak_below_limit}

\section*{S5 Origin rule}
Replacing the maximum vertex (with the plateau-centre tie rule) by the vertex with the highest mean score in its 3-mm neighbourhood gave a similar synthetic bias up to 15\,mm and a larger one above (e.g.\ $+0.50$ instead of $+0.21$\,mm at $\sigma_0 = 30$\,mm and 10\,\% noise) and moved the detected origin further from the marked pacing site in the clinical maps (median 10.1 vs 6.3\,mm; beyond 10\,mm in 23 vs 16 of 46 maps; in the narrow maps 9 vs 5 of 30); the maximum is retained.
\stab{S5}\input{tables_supp/S5_origin_rule}

\section*{S6 Non-Gaussian truth}
Fields with exponential, Lorentzian, raised-cosine and (control) Gaussian radial profiles of known full width at half maximum (7--47\,mm) were generated on the exact geodesics of the 56 meshes (17{,}920 fits: 4 profiles $\times$ 4 widths $\times$ $\eta = 0$ and 10\,\% $\times$ 56 meshes $\times$ 10 seeds) and evaluated with the Gaussian pipeline.  The Gaussian-equivalent width $2.355\,\hat\sigma$ recovered the true FWHM within $-10$ to $+4$\,\% for Gaussian and raised-cosine truth, but overestimated it by $+25$ to $+64$\,\% for exponential and by $+6$ to $+34$\,\% for Lorentzian truth at 10\,\% noise (mean error per row; Gaussian and raised-cosine truth $-8$ to $+1$\,\% at 10\,\% noise).  An exponential fit to the same binned profile had the lower AIC ($\Delta\mathrm{AIC} > 0$; the clinical counts of the manuscript use $\Delta\mathrm{AIC} > 2$) in at least 96\,\% of the exponential and Lorentzian fields, except for the Lorentzian at 47\,mm (76\,\%; Gaussian preferred in 24\,\%); a lower AIC of the exponential is a model preference among the two candidates, not an identification of a Lorentzian truth, which was not fitted.  For raised-cosine truth the Gaussian was preferred (a raised-cosine model was not fitted as a candidate).
\stab{S6}\input{tables_supp/S6_nongaussian}

\section*{S7 Anisotropic truth on the sphere}
Elliptical Gaussian fields with principal widths $\sigma_a/\sigma_b = 1.5$, 2, 3 and 4 (geometric mean 5 and 10\,mm) were evaluated with the isotropic pipeline: without noise $\hat\sigma$ equalled the geometric mean $\sqrt{\sigma_a \sigma_b}$ within 1\,\%, 1--3\,\%, 6--9\,\% and 10--14\,\%, respectively; at 10\,\% noise $\hat\sigma/\sqrt{\sigma_a \sigma_b}$ was 0.87--1.04 (isotropic control 1.02), with a median $R^2 \ge 0.97$ in every condition.
\stab{S7}\input{tables_supp/S7_anisotropy_sphere}
\begin{figure}[htbp]\centering\includegraphics[width=0.5\linewidth]{figS_anisotropy_sphere.pdf}
\caption{Isotropic width recovered from elliptical Gaussian fields on the sphere relative to the geometric mean of the principal widths, by axis ratio (median and 5th--95th percentile of 20 trials; filled: noise-free, open: 10\,\% vertex noise).}\end{figure}

\section*{S8 Anisotropy of the clinical maps (exploratory)}
\noindent\textit{Method.} Geodesic polar coordinates $(r,\varphi)$ around the detected origin were obtained from the logarithmic map of the Vector Heat Method (Sharp et al., ACM Trans. Graph. 2019), computed with the point-cloud solver of potpourri3d on the mesh vertices (the exported triangle connectivity is not used; the solver builds its own neighbourhoods and operators from the vertex positions, so boundary geometry and irregular sampling can still influence the log map; radius error against exact geodesics within 40\,mm: mean $-1.5$\,\%, 95th percentile 4\,\%). A radial Gaussian fitted to the binned log-map radii gave $\sigma_{\rm radial}$ and defined the fit window $r \le \mathrm{clip}(4\sigma_{\rm radial}, 12, 60)$\,mm. Within the window the scores were binned on an equal-area polar grid (2-mm rings; $\mathrm{round}(2\pi(k+\tfrac12))$ sectors in ring $k$, at most 32; cells with at least three vertices; equal cell weights as in the radial pipeline). An isotropic Gaussian (amplitude, width, baseline, centre $x_0,y_0$; five parameters) and an elliptical Gaussian (amplitude, two principal widths, orientation, baseline, centre; seven parameters) were fitted by bounded least squares (widths 0.1--200\,mm, centre within $\pm 1$\,mm of the origin on each axis, no re-centring), the ellipse from four starting orientations ($-90^\circ$, $-45^\circ$, $0^\circ$, $45^\circ$), keeping the fit with the smallest residual. $\Delta$AIC compares the two models (positive: ellipse preferred). The null distribution of the axis ratio was obtained by permuting the angular positions of the vertices within each 2-mm ring (100 permutations; radial profile preserved, direction destroyed) and refitting; $p = (\#\{\text{null ratio} \ge \text{observed}\}+1)/(n+1)$.

\noindent\textit{Results.} In the 33 maps with an isotropic width below 30\,mm in this fit and with both principal widths inside the fit bounds, the ratio of the principal widths had a median of 2.7 (IQR 2.0--5.3) and exceeded the within-ring permutation null in 24; in 9 further maps one principal width ran to a fit bound.  The permutation null removes all angular structure, including smooth structure that is not elliptical, and the upper quartile of the ratios lies beyond the axis ratios of the synthetic tests (Table~S7).

\stab{S8}\input{tables_supp/S8_anisotropy_clinical}

\section*{S9 Structured and correlated deviations on the clinical meshes}
On the 56 clinical meshes (exact-geodesic truth; $\sigma_0 = 5$, 10 and 15\,mm; 15{,}120 fields) the Gaussian decay was combined with two kinds of structured deviation.  (a)~Smooth random fields: 30 Gaussian bumps of 8\,mm width with an amplitude SD of 0.10--0.50\,$A$, plus 10\,\% vertex noise.  (b)~Spatially correlated noise in place of independent noise: white noise~$\varepsilon$ smoothed by one implicit diffusion step, $(M_{\mathrm{mass}} + \ell^2 L)\,u = M_{\mathrm{mass}}\,\varepsilon$, with smoothing scale $\ell$ = 5 and 10\,mm, rescaled to an SD of 0.10--0.40\,$A$.  Here $\ell$ is a smoothing parameter, not an estimated correlation length.

\textit{Scatter.}  Correlated noise of 0.25\,$A$ produced a within-bin scatter of 15--19\,\% of~$A$ (5.1--6.8 score points at $A = 35$), and 0.40\,$A$ one of 19--26\,\%.  Relative to the amplitude, the clinical maps scatter more (36\,\% of~$A$, 7.7 points).  Direction-dependent decay adds within-bin scatter of up to 20\,\% of~$A$ at an axis ratio of~4, with little effect on~$R^2$, the origin or the width (Tables~S7 and S10).

\textit{Fit quality and origin.}  With correlated noise of 0.25\,$A$ and a smoothing scale of 10\,mm, the median $R^2$ was 0.89, 13\,\% of the fits were flagged, and the detected origin was displaced by more than 10\,mm in 26\,\% of the fields (Table~S9b).  At 0.40\,$A$ these figures were 0.85, 17\,\% and 64\,\%.  At 0.40\,$A$ the origin was displaced in 46--80\,\% of the fields.  Bumps of 0.25\,$A$ lowered the median $R^2$ to 0.93--0.98 and displaced the origin in 2--6\,\% of the fields; bumps of 0.50\,$A$ displaced it in 55--65\,\% (Fig.~S3).

\textit{Width.}  Among trials with the detected origin within 10\,mm of the true source, the median width bias at 0.25\,$A$ was within $\pm 0.7$\,mm, with an SD of 1.8--3.5\,mm.  At 0.10\,$A$ the bias stayed small, but the SD of the width was larger than under independent noise (0.6--1.5 vs 0.3--0.6\,mm at $\sigma_0$ = 5--15\,mm).

\textit{Focality index.}  Independent vertex noise of 10\,\% on the meshes did not affect the focality index (no field flagged).  It flagged 45--69\,\% of the fields with correlated deviations of 0.25\,$A$ and 21--74\,\% of those with bumps of 0.25\,$A$ at $\sigma_0 = 5$--15\,mm (Table~S14b).

\textit{Comparison with the clinical maps.}  No single condition reproduces the summary behaviour of the clinical fits.  The clinical maps had a lower median $R^2$ (0.82) and a higher flag rate (24\,\%) than the two correlated-noise conditions with a 10-mm smoothing scale.  Their proportion of maxima more than 10\,mm from the marked pacing site (35\,\% of the 46 maps with such a site) is an offset from a marker, not a displacement of the origin by noise, and cannot be compared with the simulated displacement.  The conditions illustrate mechanisms; they are not a model of the clinical noise.
\stab{S9}\input{tables_supp/S9_structured_meshes}
\stab{S9b}\input{tables_supp/T1b_validation_correlated}
\begin{figure}[htbp]\centering\includegraphics[width=0.95\linewidth]{figS_structured_deviations.pdf}
\caption{Behaviour of the pipeline under structured deviations on the 56 clinical meshes: median $R^2$, fraction of fits flagged by $R^2 < 0.7$ and fraction of fields in which the detected origin lies more than 10 mm from the true one, against the amplitude of the deviation (bumps with 10\,\% independent noise; correlated noise with 5 and 10 mm smoothing scale) for $\sigma_0$ = 5, 10 and 15 mm.  Horizontal lines: clinical median $R^2$ and clinical flag rate.}\end{figure}

\section*{S10 Sources of within-bin scatter}
\stab{S10}\input{tables_supp/S10_scatter_sources}

\section*{S11 Classification flags on fields of known type}
On the sphere (2{,}562 vertices) fields of known type---Gaussian ($\sigma_0 = 5$, 10, 15\,mm), flat, ring-shaped (radius 15\,mm) and two-peaked (separation 10, 15, 25, 40\,mm; amplitude ratio 0.3--1.0)---were generated with vertex noise of 5--25\,\% and classified with the two thresholds (20 trials per condition, i.e.\ steps of 5\,\%; 60 trials for the single-peak fields pooled over the three widths).  Two-peaked fields are the pointwise maximum of two Gaussian lobes of equal width ($\sigma$ = 7\,mm), not their sum.  The results below were obtained with the current estimator (plateau-centre origin, amplitude unconstrained); $F$ is unavailable when the fitted amplitude is not positive or the half-maximum set has fewer than three vertices (one implementation, \texttt{geodesic\_decay.analysis.focality\_index}, for the study scripts and the public function; on the sphere: flat fields 20--25\,\% of the trials, all other fields none, Table~S11c; on the clinical meshes 1 of 2{,}520 single-peak fields, Table~S14b); $F$ rates are fractions of the fields with $F$ available.  $F$ rates refer to the flag as applied in the pipeline ($F < 0.70$ and $R^2 \ge 0.7$).  The $R^2$ flag rejected no Gaussian field at thresholds up to 0.8 and accepted flat fields as adequate in 0, 10, 15 and 10\,\% of the trials at 5, 10, 15 and 25\,\% noise (threshold 0.7; 0, 5, 15 and 0\,\% at 0.8); ring-shaped and two-peaked fields were not rejected ($R^2 \approx 0.87$ and $0.98$).  At 5--10\,\% noise the focality flag was raised in 90--100\,\% of the trials for two peaks 25\,mm or more apart with an amplitude ratio of 0.7 or higher ($F$ = 0.48--0.64 at 10\,\% noise), in none for a ratio of 0.5, whose second lobe reaches only about the half maximum, so that its vertices above the threshold form a small separate component ($F$ = 0.83--0.85, above the threshold of 0.70), and in 0--15\,\% for peaks 15\,mm or less apart, which merge into one component; ring-shaped fields were not flagged.  At 25\,\% noise the half-maximum set fragments: single- and two-peak fields were flagged in 90--100\,\% of the trials, ring-shaped fields in 40\,\%, flat fields (inadequate fit) in none (at 10\,\% noise in 12\,\%).  Flat fields had $F < 0.70$ in 88--94\,\% of the trials with $F$ available but almost always an inadequate fit, so that they are caught by the $R^2$ flag rather than by~$F$.  With independent vertex noise of 5, 10 and 15\,\% a single-peaked field was flagged in 0, 0 and 2\,\% of the trials (25\,\%: 90\,\%), with correlated noise of 0.10\,$A$ in none and of 0.25\,$A$ in 33\,\%.  On the 55 clinical maps, $R^2$ thresholds of 0.6, 0.7 and 0.8 flag 7, 13 and 23 maps, and the median~$\sigma$ of the remaining maps is 8.5, 8.7 and 7.2\,mm.
\stab{S11a}\input{tables_supp/S11a_r2_threshold}
\stab{S11b}\input{tables_supp/S11b_focality_detection}
\stab{S11c}\input{tables_supp/S11c_focality_other}
\begin{figure}[htbp]\centering\includegraphics[width=0.95\linewidth]{figS_focality_synthetic.pdf}
\caption{Detection of a second peak by the focality index on the sphere at 10\,\% vertex noise: fraction of fields flagged by peak separation and amplitude ratio for $F$ (left; $F < 0.70$ with $R^2 \ge 0.7$ as in the pipeline, $F$ available in all two-peak fields; as Table~S11b) and the percentile index $F_{\rm pct}$ ($F_{\rm pct} < 0.70$; 95th-percentile set, largest component; right).}\end{figure}

\section*{S12 Per-map results}
In the 13 maps flagged by $R^2 < 0.7$ the bin means beyond $3\sigma$ varied with an SD of 3.8 score points (1.8 in the unflagged maps).
\stab{S12}\input{tables_supp/S12_per_map_long}

\section*{S13 Distance definitions}
Recomputing the width from the same origin with exact polyhedral geodesics changed~$\sigma$ by a median of $+0.08$\,mm (IQR $-0.03$ to $+0.20$\,mm; $\rho = 0.97$; $|\Delta\sigma| > 1$\,mm in 7 of 55 maps and in 3 of the 35 maps with $R^2 \ge 0.7$ and $\sigma < 30$\,mm); Dijkstra distances on the edge graph, which overestimate the geodesic by 5.7\,\% within 30\,mm of the origin, widened it by a median of $+0.49$\,mm (IQR $+0.23$ to $+1.09$\,mm; $> 1$\,mm in 16 maps); the Euclidean distance, shorter than the geodesic by 3.8\,\% within 30\,mm and by 15\,\% over the whole mesh (medians over maps of the mean relative error; medians of the within-map 95th percentiles 11 and 38\,\%), changed it by a median of $-0.29$\,mm (IQR $-1.51$ to $+0.14$\,mm; $\rho = 0.94$), but by more than 3\,mm in 16 broad or folded maps (8 of the 35 maps with $R^2 \ge 0.7$ and $\sigma < 30$\,mm) and with a poorer fit (median $R^2$ 0.78; 17 flagged maps).  Group medians changed little (double-chamber 5.2 Heat Method, 5.5 exact, 5.8 Dijkstra, 5.5\,mm Euclidean; single-chamber 10.4, 11.0, 10.7 and 10.8\,mm).  Exact geodesics took 0.16\,s per map (median; 1.2\,s on the largest mesh), the Heat Method 0.03\,s.
\stab{S13}\input{tables_supp/S13_distance_definitions}

\section*{S14 Focality index of the clinical maps}
Two of the 13 maps with an inadequate fit ($R^2 < 0.7$) also had $F < 0.70$; there $F$ is not interpreted (Fig.~6B of the main text).
\stab{S14}\input{tables_supp/S14_focality_clinical}
\stab{S14b}\input{tables_supp/S14b_focality_meshes}

\section*{S15 Mixed model}
Model and diagnostics.  The linear mixed model of $\log\sigma$ (reference and region as fixed effects, patient random intercept) was fitted by REML.  The default optimiser of statsmodels did not converge for the 55 maps (gradient 0.076 at a random-intercept variance of 0.0003); the REML log-likelihood profiled over the variance ratio $\gamma = \tau^2/\sigma^2$ in closed form has its maximum at $\gamma = 0.028$ (patient ICC 0.03) and is flat near zero (difference 0.03 log-likelihood units between $\gamma = 0$ and the optimum).  All estimates of the manuscript and of Table~S15 are taken at this profiled optimum, with the model-based covariance (inverse GLS information at the profiled variance ratio) and an approximate $t$ reference with (patients $-$ 2) degrees of freedom, applied to all terms.  This reference is motivated by the reference configuration being a between-patient factor with 8 and 8 patients; it is not a Kenward--Roger or Satterthwaite adjustment, does not propagate the uncertainty of the variance component, and its calibration was not verified.  The ML fits of the full and reduced models used for the likelihood-ratio test converged and agree with the profiled ML maxima ($\chi^2 = 5.21$, $p = 0.022$).  Small-sample alternatives that do not rely on the random effect gave similar results: ordinary least squares with CR2 patient-cluster-robust standard errors and Satterthwaite degrees of freedom 1.79 (1.04--3.10), $p = 0.039$; a wild cluster bootstrap (Webb weights, 9{,}999 draws) $p = 0.031$.  A two-stage patient-level sensitivity analysis (residuals of $\log\sigma$ on region, averaged per patient, Welch test, 8 vs 8; $\exp$ of the difference of the group means; the uncertainty of the first stage is not propagated) gave 1.75 (0.93--3.29), $p = 0.08$, and the patient medians were 11.2 vs 5.0\,mm (Mann--Whitney $p = 0.33$).  With widths above 30\,mm clipped at 30\,mm the ratio was 1.70 (1.05--2.74); restricted to the 42 maps with $R^2 \ge 0.7$ (13 patients: 6 single, 7 double) it was 2.36 (1.21--4.61).  The ratio of the map-level group medians was 1.97 (patient-cluster bootstrap 95\,\% interval 1.35--2.94).  Widths in the broad and diffuse ranges enter these models with their numerical value; the clipped analysis is a sensitivity, not a censored-data model.  The estimates describe an association with low precision, not a confirmed effect; the several analyses of the same 16 patients are not independent confirmations.
\stab{S15}\input{tables_supp/S15_mixed_model}

\section*{S16 Runtime}
The internal factor of the Heat-Method library is not exposed.  As a proxy for its sparsity we factorised the heat-flow matrix of Section 3.3 of the manuscript, written $M_{\mathrm{mass}} + t\,L$ with the positive semidefinite Laplacian $L = -L_C$, of the 59 clinical meshes with SuperLU in symmetric mode (minimum-degree ordering): the matrix has 7 non-zeros per row, its factor a median of 18 (range 12--36), i.e.\ a fill-in of about 4.4 relative to the lower triangle of the matrix (about 4 non-zeros per row), and the number of non-zeros of the factor scaled as $N^{1.28}$.  The $\mathcal{O}(N^{1.5})$ bound of the main text holds for nested-dissection orderings; the measurements refer to the library's default ordering.  Wall times were measured on a single core of a Linux virtual machine on Apple silicon, as the minimum of three solver constructions and of six queries from two source vertices; the log-log fit of the construction time against $N$ had $r = 0.996$.
\stab{S16}\input{tables_supp/S16_runtime_long}
\begin{figure}[htbp]\centering\includegraphics[width=0.5\linewidth]{figS_runtime.pdf}
\caption{Wall time of solver construction and of a single-source query against the number of vertices of the 59 clinical meshes (single core), with log-log fits.}\end{figure}

\section*{S17 Bin width and mesh smoothing}
Taubin smoothing of the meshes (10 and 30 iterations) changed the widths by more than 1\,mm only in one broad map with a large plateau of tied maximum scores (36 vs 41--42\,mm).  The Euclidean chord is a lower bound on the geodesic distance, but not a bound on the fitted width.
\stab{S17}\input{tables_supp/S17_binwidth_smoothing}

\section*{S18 Marked pacing site as a reference for the origin}
Using the marked site instead of the maximum as origin gave similar median widths when the origins were within 5\,mm (19 maps: median difference $+0.18$\,mm, IQR $-0.37$ to $+0.82$\,mm).  When they were more than 10\,mm apart, the median difference was $+9.3$\,mm (16 maps; two fits reached the parameter bound).
\stab{S18}\input{tables_supp/S18_marker_check}
The distances of the 46 maps against the width are shown in Fig.~7 of the main text.  Among the 30 maps with $\sigma < 10$\,mm (13 patients), the maximum lay within 10\,mm of the marked site in 17 of the 21 maps with $R^2 \ge 0.7$ and in 14 of the 17 maps without either flag; the five narrow maps with a larger distance (maps 3, 9, 12, 31 and 53) lay 20.3--26.9\,mm away.
\section*{S19 Ties of the maximum score}
Rule.  When several vertices share the maximum score, the tied vertices are grouped into connected regions on the mesh edge graph, the largest region is taken (on equal size, the one containing the lowest vertex index), and the origin is its vertex nearest to the centroid of that region (plateau centre).
The maximum score was shared by several vertices in 14 maps (2 to 7{,}369 tied vertices).  With the first tied vertex in vertex order as origin instead of the plateau centre, $\sigma$ changed by more than 1\,mm in 7 of them (3 with an adequate fit and $\sigma < 30$\,mm, by 1.3--3.1\,mm), the $R^2$ flag changed in 4 (in both directions, so that 13 maps remained flagged) and the width class in 3; the median widths of the cohort and of both reference groups were unchanged.
\stab{S19}\input{tables_supp/S19_argmax_ties}
\section*{S20 Gaussian versus exponential fit of the clinical profiles}
Over the fitted range (median 107\,mm), the exponential fit had a median $R^2$ of 0.84 against 0.82 for the Gaussian, and a higher $R^2$ in 33 maps.  The median root-mean-square residual was 2.3 score points for the exponential and 2.4 for the Gaussian.  Within $3\hat\sigma$ of the origin both models had a median residual of 1.9 points, suggesting that the differences lie mainly in the far field.  In a generalised-Gaussian fit the exponent~$\beta$ was at a bound of the fit (0.3 or 4.0) in 12 maps and within 0.01 of it in one more.

\textit{Empirical half-maximum width.}  The baseline is the median of the bin means beyond $\max(3\hat\sigma, 6\,\mathrm{mm})$, or the lowest bin mean if fewer than three such bins exist.  The empirical half-maximum width is twice the centre of the first bin whose mean falls below half the height of the highest bin mean above this baseline.  It uses $\hat\sigma$ only for the baseline window and, because the bin centres lie at odd millimetres, takes values on a 4-mm grid.  On the clinical maps it was a median of $1.00 \times 2.355\,\hat\sigma$ (IQR 0.83--1.13).
\stab{S20}\input{tables_supp/S20_model_comparison}
\section*{S21 Width by region and reference configuration}
\begin{figure}[htbp]\centering\includegraphics[width=0.5\linewidth]{figS_boxplot_sigma.pdf}
\caption{Width~$\sigma$ of the 55 pace maps by anatomic region and reference-catheter configuration (points: maps; boxes: median and interquartile range; diamonds: median per patient and region; logarithmic axis; dashed: 10\,mm).  Exploratory; mixed-model estimates in Table~2 of the main text and Section~S15.}\end{figure}
\end{document}
